% GENERATED FILE. Edit canonical lessons, then run npm run generate:books.
% canonical-source-hash: fnv1a64:879e449d21ac504b
% canonical-lessons: IT-C424-boa, IT-C424-approdare, IT-C424-ormeggiare, IT-C424-salvagente, IT-C424-nausea

\chapter{Arrivare in porto}
\label{ch:it-arrivare-in-porto}
\hlchaptermodality{424}

\begin{chapteropening}
\textbf{By the end of this chapter:} \emph{I can talk about reaching land: the buoy, coming ashore, mooring, the life jacket and feeling sick.}

Complete the last lesson of chapter 424: I can talk about reaching land: the buoy, coming ashore, mooring, the life jacket and feeling sick.
\end{chapteropening}

\section[la boa]{la boa — the buoy}

\label{lesson:IT-C424-boa}



\begin{quote}
Before the new one: say the Italian for the bow, then the Italian for the stern.
\end{quote}

\subsection*{You'll want to know: la boa}
\textbf{la boa} — \textquotedblleft{}the buoy\textquotedblright{}.

\begin{quote}
Non nuotare oltre le boe.
\end{quote}

\begin{quote}
\emph{Do not swim beyond the buoys.}
\end{quote}

On beaches the buoys mark the safe swimming area.

\begin{cousinweb}
\emph{boa} is usually traced to a Germanic word for a sign or marker, through Old French, but the origin is not settled.
\end{cousinweb}

\begin{grammarlens}[title={What you've built}]
One more word for telling, explaining and describing at B1.
\end{grammarlens}

\subsection*{Your turn}
\begin{itemize}
\raggedright
  \item \emph{Say it:} \emph{la boa}
  \item \emph{Say it:} \emph{la boa}, once more
  \item \emph{Say it:} a sentence of your own with \emph{la boa}
  \item \emph{Recall:} say the Italian for the bow, then the Italian for the stern, then say \emph{la boa} again
\end{itemize}

\begin{tcolorbox}[breakable,colback=teal!4,colframe=teal!35!black,title={Before you move on}]
What does la boa mean? (\textquotedblleft{}The buoy\textquotedblright{}.) Say it once more.
\end{tcolorbox}
\section[approdare]{approdare — to come ashore, to land (a boat)}

\label{lesson:IT-C424-approdare}



\begin{quote}
Before the new one: say the Italian for the stern, then the Italian for the buoy.
\end{quote}

\subsection*{You'll want to know: approdare}
\textbf{approdare} — \textquotedblleft{}to come ashore, to land (a boat)\textquotedblright{}.

\begin{quote}
Siamo approdati su una piccola isola.
\end{quote}

\begin{quote}
\emph{We came ashore on a small island.}
\end{quote}

Figuratively: \emph{approdare a un'idea} — to arrive at an idea.

\begin{cousinweb}
\emph{approdare} is built from \emph{a-} and \emph{proda} 'shore, edge', which is usually traced to Latin \emph{prōra} 'prow' (Lewis \& Short), an inherited word.
\end{cousinweb}

\begin{grammarlens}[title={What you've built}]
One more word for telling, explaining and describing at B1.
\end{grammarlens}

\subsection*{Your turn}
\begin{itemize}
\raggedright
  \item \emph{Say it:} \emph{approdare}
  \item \emph{Say it:} \emph{approdare}, once more
  \item \emph{Say it:} a sentence of your own with \emph{approdare}
  \item \emph{Recall:} say the Italian for the stern, then the Italian for the buoy, then say \emph{approdare} again
\end{itemize}

\begin{tcolorbox}[breakable,colback=teal!4,colframe=teal!35!black,title={Before you move on}]
What does approdare mean? (\textquotedblleft{}To come ashore, to land (a boat)\textquotedblright{}.) Say it once more.
\end{tcolorbox}
\section[ormeggiare]{ormeggiare — to moor}

\label{lesson:IT-C424-ormeggiare}



\begin{quote}
Before the new one: say the Italian for the buoy, then the Italian for to come ashore.
\end{quote}

\subsection*{You'll want to know: ormeggiare}
\textbf{ormeggiare} — \textquotedblleft{}to moor\textquotedblright{}.

\begin{quote}
Ormeggiamo la barca vicino al molo.
\end{quote}

\begin{quote}
\emph{We moor the boat near the pier.}
\end{quote}

\emph{L'ormeggio} — the mooring.

\begin{cousinweb}
\emph{ormeggiare} is usually traced to Greek \emph{hórmos} 'anchorage, roadstead', but the path into Italian is not settled.
\end{cousinweb}

\begin{grammarlens}[title={What you've built}]
One more word for telling, explaining and describing at B1.
\end{grammarlens}

\subsection*{Your turn}
\begin{itemize}
\raggedright
  \item \emph{Say it:} \emph{ormeggiare}
  \item \emph{Say it:} \emph{ormeggiare}, once more
  \item \emph{Say it:} a sentence of your own with \emph{ormeggiare}
  \item \emph{Recall:} say the Italian for the buoy, then the Italian for to come ashore, then say \emph{ormeggiare} again
\end{itemize}

\begin{tcolorbox}[breakable,colback=teal!4,colframe=teal!35!black,title={Before you move on}]
What does ormeggiare mean? (\textquotedblleft{}To moor\textquotedblright{}.) Say it once more.
\end{tcolorbox}
\section[il salvagente]{il salvagente — the lifebelt, the life jacket}

\label{lesson:IT-C424-salvagente}



\begin{quote}
Before the new one: say the Italian for to come ashore, then the Italian for to moor.
\end{quote}

\subsection*{You'll want to know: il salvagente}
\textbf{il salvagente} — \textquotedblleft{}the lifebelt, the life jacket\textquotedblright{}.

\begin{quote}
Sulla nave ogni posto ha un salvagente.
\end{quote}

\begin{quote}
\emph{On the ship every seat has a life jacket.}
\end{quote}

Built from \emph{salva} and \emph{gente}: it saves people. Plural unchanged.

\begin{cousinweb}
\emph{salvagente} is built from \emph{salvare} 'to save', inherited from Late Latin \emph{salvāre} (EtymDB), and \emph{gente} 'people', inherited from Latin \emph{gēns}.
\end{cousinweb}

\begin{grammarlens}[title={What you've built}]
One more word for telling, explaining and describing at B1.
\end{grammarlens}

\subsection*{Your turn}
\begin{itemize}
\raggedright
  \item \emph{Say it:} \emph{il salvagente}
  \item \emph{Say it:} \emph{il salvagente}, once more
  \item \emph{Say it:} a sentence of your own with \emph{il salvagente}
  \item \emph{Recall:} say the Italian for to come ashore, then the Italian for to moor, then say \emph{il salvagente} again
\end{itemize}

\begin{tcolorbox}[breakable,colback=teal!4,colframe=teal!35!black,title={Before you move on}]
What does il salvagente mean? (\textquotedblleft{}The lifebelt, the life jacket\textquotedblright{}.) Say it once more.
\end{tcolorbox}
\section[la nausea]{la nausea — the nausea, the feeling sick}

\label{lesson:IT-C424-nausea}



\begin{quote}
Before the new one: say the Italian for to moor, then the Italian for the lifebelt.
\end{quote}

\subsection*{You'll want to know: la nausea}
\textbf{la nausea} — \textquotedblleft{}the nausea, the feeling sick\textquotedblright{}.

\begin{quote}
Con le onde ho avuto la nausea.
\end{quote}

\begin{quote}
\emph{With the waves I felt sick.}
\end{quote}

\emph{Avere la nausea}; \emph{il mal di mare} is seasickness.

\begin{cousinweb}
\emph{nausea} is a learned borrowing from Latin \emph{nausea}, from Greek \emph{nausía} 'seasickness', built on \emph{naûs} 'ship' (EtymDB).
\end{cousinweb}

\begin{grammarlens}[title={What you've built}]
One more word for telling, explaining and describing at B1.
\end{grammarlens}

\subsection*{Your turn}
\begin{itemize}
\raggedright
  \item \emph{Say it:} \emph{la nausea}
  \item \emph{Say it:} \emph{la nausea}, once more
  \item \emph{Say it:} a sentence of your own with \emph{la nausea}
  \item \emph{Recall:} say the Italian for to moor, then the Italian for the lifebelt, then say \emph{la nausea} again
\end{itemize}

\begin{tcolorbox}[breakable,colback=teal!4,colframe=teal!35!black,title={Before you move on}]
What does la nausea mean? (\textquotedblleft{}The nausea, the feeling sick\textquotedblright{}.) Say it once more.
\end{tcolorbox}
